\scanpage{47}
\mspage{2}
% archive foliation: 46
\[
(5)\qquad\sigma_3(z)=-\frac1z\qquad
\Delta_3=V(z^2+1)\quad(=\{i,-i\}\text{ si }\exists\text{ racine primitive de }\unsure{\epsilon-1|4\text{-}1}\text{ sur }S)
\]
On trouve donc

\uline{Proposition} Loc pour la top. étale, une représentation monomorphique de $\mathbb D_4=\mathbb F_2\times\mathbb F_2$ dans $\operatorname{Aut}(\mathbb P^1_S)$ (sur une base $S$ à car. résid. $\ne2$) est donnée par (4) (où $\sigma_1,\sigma_2$ sont les autos correspondant aux éléments de la base canonique $(e_1,e_2)$ de $\mathbb D_4$ sur $\mathbb F_2$) - qui donne bien une représentation monomorphique.

Posons
\[
(6)\qquad\Delta=\Delta_1\amalg\Delta_2\amalg\Delta_3
\qquad(=\{0,\infty,1,-1,i,-i\}\text{ dans le cas normalisé})
\]
on appelle $\Delta$ la \uline{configuration octaédrale} associée à l'opération de $V$ sur \deleted{$\unsure{\leftarrow}$} $X$ - par descente plate, elle est définie (comme diviseur étale de degré $6$ de $X/S$) dès qu'on a une opération monomorphique \inserted{(dans chaque fibre)} d'un groupe $V$ \inserted{$\simeq\unsure{\mathbb D_3|\mathbb D_2}$} sur $X$, ou plus généralement d'un schéma en groupes $\underline V$ sur $S$ loc. isom. au schéma en groupes constant. \deleted{\unsure{Alors}} De façon générale, on appelle \uline{configuration octaédrale} dans $X/S$ un diviseur étale de degré $6$ de $X/S$, qui soit
